% Source: https://artofproblemsolving.com/community/c3h2549833_arithmetic_math_2
% Stable identity: aops-topic-2549833-q01
\emph{All in the same direction.} If all six dogs run clockwise (or all counterclockwise), at every moment each has covered the same distance along the boundary, so the distances between them along the boundary never change, and no two meet.

\emph{Otherwise.} Going around the hexagon, some dog runs clockwise while the next dog in the clockwise direction runs counterclockwise. These two start at the two ends of the same side, of length $L$, and run towards each other along it. Unless a collision has already happened somewhere, they meet at the midpoint of that side after time $\dfrac{L}{2v}$, where $v$ is the speed. Either way there is a collision.

So no collision happens exactly when all dogs choose the same direction, with probability
\[\dfrac{2}{2^6}=\boxed{\dfrac1{32}}.\]
The side lengths do not matter.
